\exercisegroup

\begin{enumerate}
\def\labelenumi{\arabic{enumi})}
\tightlist
\item
  设 \((\mathrm{Q}_{i})_{i\in\mathrm{I}}\) 是一族箭图。用 S 表示族 \((\mathrm{Som}(\mathrm{Q}_{i}))\) 的和集，用 F 表示族 \((\mathrm{Fl}(\mathrm{Q}_{i}))\) 的和集，并用 o 与 t 表示由诸箭图 \(Q_{i}\) 的起点映射与终点映射诱导的从 F 到 S 的映射。
\end{enumerate}

\begin{enumerate}
\def\labelenumi{\alph{enumi})}
\item
  证明 \(\mathrm{Q} = (\mathrm{S}, \mathrm{F}, o, t)\) 是一个箭图。称它为族 \((\mathrm{Q}_{i})\) 的和箭图。
\item
  对 \(i \in I\)，存在唯一的箭图态射 \(j_{i}\)（由 \(Q_{i}\) 到 Q），使得映射 \(\operatorname{Som}(j_{i})\) 是从 \(\operatorname{Som}(Q_{i})\) 到 S 中的典范映射，且映射 \(\operatorname{Fl}(j_{i})\) 是从 \(\operatorname{Fl}(Q_{i})\) 到 S 中的典范映射。
\item
  设 K 是一个箭图，且对每个 \(i \in I\)，设 \(\varphi_{i}\) 是从 \(Q_{i}\) 到 K 的箭图态射。证明存在唯一的从 Q 到 K 的箭图态射 \(\varphi\)，使得 \(\varphi \circ j_{i} = \varphi_{i}\) 对所有 \(i \in I\) 成立。
\end{enumerate}

\begin{enumerate}
\def\labelenumi{\arabic{enumi})}
\setcounter{enumi}{1}
\tightlist
\item
  设 \(Q = (S, F, o, t)\) 是一个箭图。对每个 \(s \in S\)，称 \(d_{+}(s)\)（相应地，\(d_{-}(s)\)）为 Q 在 s 处的出次数（相应地，入次数），并把 \(d_{+}(s)\)（相应地，\(d_{-}(s)\)）定义为集合 \(\{f \in F \mid o(f) = s\}\)（相应地，\(\{f \in F \mid t(f) = s\}\)）的基数。若 \(d_{+}(s)\) 与 \(d_{-}(s)\) 对每个 \(s \in S\) 均有限，则称 Q 是局部有限的；若集合 F 与 S 均有限，则称 Q 是有限的。
\end{enumerate}

证明等式 \(\sum_{s\in S}d_{+}(s) = \sum_{s\in S}d_{-}(s) = \mathrm{Card}(\mathrm{F})\)

\begin{enumerate}
\def\labelenumi{\arabic{enumi})}
\setcounter{enumi}{2}
\tightlist
\item
  设 \(\mathrm{Q} = (\mathrm{S},\mathrm{F},o,t)\) 是一个局部有限箭图。我们定义矩阵 \(M_{\mathrm{Q}} = (m_{i,j})\)（类型为 (S,S)，元素在 \(\mathbf{Z}\) 中）如下：
\end{enumerate}

\[
 m _ {i, j} = \operatorname{Card} \left\{f \in \mathrm{F} \mid o (f) = i \text {and} t (f) = j \right\}.
\]

矩阵 \(M_{\mathrm{Q}}\) 称为 Q 的邻接矩阵。

\begin{enumerate}
\def\labelenumi{\alph{enumi})}
\item
  给定箭图 Q 与 Q'，给出关于它们的邻接矩阵的一个必要充分条件，使 Q 与 Q' 同构。
\item
  用邻接矩阵 \(M_{Q}\) 表出 Q 中起点与终点给定的长为 n 的道路的数目。
\item
  证明一个图的底箭图的邻接矩阵是对称的。
\end{enumerate}

\begin{enumerate}
\def\labelenumi{\arabic{enumi})}
\setcounter{enumi}{3}
\tightlist
\item
  设 \(Q = (S, F, o, t)\) 是一个箭图。用 I 表示族 \((S, F)\) 的和集；把 S 与 F 等同于 I 的子集。设 k 是一个交换环。把箭图 Q 的 k-代数称为商 \(A_{Q}\)（即集合 I 上的自由结合 k-代数对由元素 \(X_{s}^{2} - X_{s}\)（对 \(s \in S\)）、\(X_{u}X_{v} - X_{v}X_{u}\)（对 \(u, v \in S\)）、\(X_{f}X_{t(f)} - X_{f}\) 与 \(X_{o(f)}X_{f} - X_{f}\)（对 \(f \in F\)）所生成的双边理想所作的商）。对 \(i \in I\)，用 \(x_{i}\) 表示 \(X_{i}\) 在代数 \(A_{Q}\) 中的像。
\end{enumerate}

\begin{enumerate}
\def\labelenumi{\alph{enumi})}
\item
  证明有 \(\sum_{s\in S}x_s = 1\)
\item
  证明若 f 与 g 是 Q 中不可合成的箭，则有 \(x_{f}x_{g}=0\)。
\item
  对 Q 中每条道路 \(c = (a_{0}, f_{1}, \ldots, f_{n}, a_{n})\)，令 \(x_{c} = x_{a_{0}} x_{f_{1}} \ldots x_{f_{n}}\)。设 c 与 \(c'\) 是 Q 中的道路；证明有 \(x_{c} x_{c'} = x_{cc'}\)（若 c 与 \(c'\) 可合成），否则有 \(x_{c} = x_{c'}\)。
\item
  设 C 是箭图 Q 中道路的集合；证明族 \((x_{c})_{c\in C}\) 是 k-模 \(A_{Q}\) 的一组基。
\item
  设 \(J_{Q}\) 是 \(A_{Q}\) 中由形如 \(x_{f}\)（\(f \in F\)）的元素生成的双边理想。证明代数 \(A_{Q}/J_{Q}\) 同构于 \(k^{(S)}\)。
\item
  为使 k-代数 \(A_{Q}\) 是一个有限生成的 k-模，必要且充分的是 Q 的顶点集有限且 Q 中每个圈的长为零。此时理想 \(J_{Q}\) 是幂零的。
\end{enumerate}

在 §1 的其余习题中，我们沿用这些记号。当我们说到 \(A_{Q}\)-模（不作另外说明时），指的是右模。若 A 是一个环，A-模 M（不作另外说明时）总指右 A-模；若 a 是 A 的一个元素，用 \(a_{M}\) 或 \((a)_{M}\) 表示 M 的位似 \(x \mapsto xa\)。

\begin{enumerate}
\def\labelenumi{\arabic{enumi})}
\setcounter{enumi}{4}
\tightlist
\item
  \begin{enumerate}
  \def\labelenumii{\alph{enumii})}
  \tightlist
  \item
    证明四元组 \(Q^{\circ} = (S, F, t, o)\) 是一个箭图。称它为 Q 的反箭图。
  \end{enumerate}
\end{enumerate}

\begin{enumerate}
\def\labelenumi{\alph{enumi})}
\setcounter{enumi}{1}
\tightlist
\item
  代数 \(A_{Q^{\circ}}\) 同构于代数 \(A_{Q}\) 的反代数。
\end{enumerate}

\begin{enumerate}
\def\labelenumi{\arabic{enumi})}
\setcounter{enumi}{5}
\tightlist
\item
  \begin{enumerate}
  \def\labelenumii{\alph{enumii})}
  \tightlist
  \item
    设 \((M_s)_{s\in\operatorname{Som}(Q)}\) 是一族 \(k\)-模，\((u_f)_{f\in\operatorname{Fl}(Q)}\) 是一个族，其中对每个 \(f\in\operatorname{Fl}(Q)\)，\(u_f\) 是从 \(M_{o(f)}\) 到 \(M_{t(f)}\) 的线性映射。设 M 是族 \((M_s)\) 的直和；对 \(s\in S\)，用 \(j_s\) 表示 \(M_s\) 到 M 中的典范单射，用 \(p_s\) 表示 M 到 \(M_s\) 上的典范投影。证明在 M 上存在唯一的 \(A_Q\)-模结构，使得有 \((x_s)_M = j_s \circ p_s\) 且 \((x_f)_M = j_{t(f)} \circ u_f \circ p_{o(f)}\)（对所有 \(s \in \operatorname{Som}(Q)\) 及每个 \(f \in \operatorname{Fl}(Q)\)）。
  \end{enumerate}
\end{enumerate}

反之，每个 \(A_{Q}\)-模都具有此形式。

\begin{enumerate}
\def\labelenumi{\alph{enumi})}
\setcounter{enumi}{1}
\tightlist
\item
  若 \(\Lambda\) 是一个 \(k\)-模且 \(s \in \operatorname{Som}(\mathbf{Q})\)，用 \(\Lambda(s)\) 表示如下的 \(\mathrm{A}_{\mathrm{Q}}\)-模：它与族 \((\mathrm{M}_i)_{i \in \operatorname{Som}(\mathbf{Q})}\)（一个 \(k\)-模族，由 \(\mathrm{M}_i = \Lambda\)（若 \(i = s\)）、否则 \(\mathrm{M}_i = 0\) 给出）以及线性映射族 \((u_f)_{f \in \operatorname{Fl}(\mathbf{Q})}\)（其中 \(u_f = 0\) 对一切 \(f \in \operatorname{Fl}(\mathbf{Q})\)）相伴。
\end{enumerate}

设 \(\Lambda\) 与 \(\Lambda'\) 是非零 \(k\)-模。设 \(s\) 与 \(s'\) 是 Q 的顶点。证明 \(\Lambda(s)\) 同构于 \(\Lambda'(s')\) 当且仅当 \(\Lambda\) 与 \(\Lambda'\) 作为 \(k\)-模同构且 \(s = s'\)。

\begin{enumerate}
\def\labelenumi{\alph{enumi})}
\setcounter{enumi}{2}
\item
  若 \(\Lambda\) 是一个单 \(k\)-模，则 \(\Lambda(s)\) 都是单 \(\mathrm{A}_{\mathrm{Q}}\)-模（对每个顶点 \(s\)）。
\item
  进一步假设 Q 中每个圈的长为零。证明每个单 \(A_Q\)-模都具有形式 \(\Lambda(s)\)，其中 \(\Lambda\) 是一个单 \(k\)-模而 s 是 Q 的一个顶点。
\end{enumerate}

\begin{enumerate}
\def\labelenumi{\arabic{enumi})}
\setcounter{enumi}{6}
\tightlist
\item
  设 \(k\) 是一个代数闭域。对每个 \(s \in \operatorname{Som}(Q)\)，令 \(\mathrm{P}(s) = x_s\mathrm{A}_{\mathrm{Q}}\)。
\end{enumerate}

\begin{enumerate}
\def\labelenumi{\alph{enumi})}
\item
  证明 \(\mathrm{P}(s)\) 是一个投射且不可分解的 \(A_{Q}\)-模，且其底座同构于 \(A_{Q}\)-模 \(k(s)\)。再证明 \(\mathrm{P}(s)/\mathrm{P}(s)\mathrm{J}_{\mathrm{Q}}\) 同构于 \(k(s)\)。
\item
  假设 Q 中每个圈的长为零。证明对每个 \(s \in \operatorname{Som}(Q)\)，从 \(k\) 到 \(\operatorname{End}(\mathrm{P}(s))\) 的典范同态是一个同构。证明对每个投射且不可分解的 \(A_Q\)-模 P，存在 Q 的唯一顶点 s 使 P 同构于 \(\mathrm{P}(s)\)。
\end{enumerate}

\begin{enumerate}
\def\labelenumi{\arabic{enumi})}
\setcounter{enumi}{7}
\tightlist
\item
  对 \(s \in \operatorname{Som}(Q)\)，令 \(\mathrm{P}(s) = x_s\mathrm{A}_{\mathrm{Q}}\)。设 M 是一个 \(\mathrm{A}_{\mathrm{Q}}\)-模。
\end{enumerate}

\begin{enumerate}
\def\labelenumi{\alph{enumi})}
\tightlist
\item
  设 \(f \in \mathrm{Fl}(\mathbf{Q})\)，令 \(u = o(f)\)，\(v = t(f)\)。证明存在唯一的一对 \((\varphi_f', \varphi_f'')\)（\(\mathrm{A}_{\mathrm{Q}}\)-模同态）
\end{enumerate}

\[
\varphi_ {f} ^ {\prime} \colon \mathrm{M} x _ {u} \otimes_ {k} \mathrm{P} (v) \to \mathrm{M} x _ {u} \otimes_ {k} \mathrm{P} (u), \quad \varphi_ {f} ^ {\prime \prime} \colon \mathrm{M} x _ {u} \otimes_ {k} \mathrm{P} (v) \to \mathrm{M} x _ {v} \otimes_ {k} \mathrm{P} (v),
\]

 使得 \(\varphi_f'(m \otimes p) = m \otimes x_f p\) 且 \(\varphi_f''(m \otimes p) = mx_f \otimes p\)（对一切 \(m \in \mathrm{M}x_u\) 及每个 \(p \in \mathrm{P}(v)\)）。

\begin{enumerate}
\def\labelenumi{\alph{enumi})}
\setcounter{enumi}{1}
\item
  设 \(s \in \text{Som}(Q)\)。证明存在唯一的 \(A_{Q}\)-模同态 \(\psi_{s}: Mx_{s} \otimes_{k} P(s) \to M\)，使得 \(\psi_{s}(m \otimes p) = mp\)（对 \(p \in P(s)\) 与 \(m \in Mx_{s}\)）。
\item
  令 \(\varphi = \bigoplus \varphi_f' - \bigoplus \varphi_f''\)，\(\psi = \sum \psi_s\)。证明有 \(\mathrm{A}_{\mathrm{Q}}\)-模的正合列
\end{enumerate}

\[
0 \to \bigoplus_ {f \in \operatorname{Fl} (\mathrm{Q})} \operatorname{M} x _ {o (f)} \otimes_ {k} \operatorname{P} (t (f)) \xrightarrow {\varphi} \bigoplus_ {s \in \operatorname{Som} (\mathrm{Q})} \operatorname{M} x _ {s} \otimes_ {k} \operatorname{P} (s) \xrightarrow {\psi} \operatorname{M} \to 0.
\]

\(d\)) 证明代数 \(\mathrm{A_Q}\) 的每个右理想都是投射的。

\begin{enumerate}
\def\labelenumi{\arabic{enumi})}
\setcounter{enumi}{8}
\tightlist
\item
  \begin{enumerate}
  \def\labelenumii{\alph{enumii})}
  \tightlist
  \item
    设 M 与 N 是 \(A_Q\)-模；对 \(s \in \operatorname{Som}(Q)\)，令 \(M_s=Mx_s\)，\(N_s=Nx_s\)。对 \(\varphi = (\varphi_s)_{s \in \operatorname{Som}(Q)}\) 与 \(f \in \operatorname{Fl}(Q)\)，如下定义映射 \(\theta_f(\varphi): Mx_{o(f)} \to Nx_{t(f)}\)
  \end{enumerate}
\end{enumerate}

\[
\theta_ {f} (\varphi) (m) = \varphi_ {o (f)} (m) x _ {f} - \varphi_ {t (f)} (m x _ {f}).
\]

证明映射

\[
\theta \colon \prod_ {s \in \operatorname{Som} (\mathrm{Q})} \operatorname{Hom} _ {k} \left(\mathrm{M} _ {s}, \mathrm{N} _ {s}\right)\rightarrow \prod_ {f \in \operatorname{Fl} (\mathrm{Q})} \operatorname{Hom} \left(\mathrm{M} _ {o (f)}, \mathrm{N} _ {t (f)}\right)
\]

 \(\varphi\mapsto(\theta_{f}(\varphi))_{f\in\mathrm{Fl}(Q)}\) 是线性的。

\begin{enumerate}
\def\labelenumi{\alph{enumi})}
\setcounter{enumi}{1}
\item
  证明 \(\operatorname{Ker}(\theta)\) 同构于 \(\operatorname{Hom}_{\mathrm{A}_{\mathrm{Q}}}(\mathrm{M},\mathrm{N})\)，且 \(\operatorname{Coker}(\theta)\) 同构于 \(\operatorname{Ext}_{\mathrm{A}_{\mathrm{Q}}}^{1}(\mathrm{M},\mathrm{N})\)。证明 \(\operatorname{Ext}_{\mathrm{A}_{\mathrm{Q}}}^{p}(\mathrm{M},\mathrm{N}) = 0\)（对每个整数 \(p \geqslant 2\)）。
\item
  设 k 是一个交换域且箭图 Q 是有限的。若 M 与 N 是有限维 k-向量空间，则 \(\operatorname{Ext}_{\mathrm{A}_{\mathrm{Q}}}^{p}(\mathrm{M},\mathrm{N})\) 对每个整数 \(p \geqslant 0\) 是一个有限维 k-向量空间，且当 \(p \geqslant 2\) 时为零，并有
\end{enumerate}

\[
\begin{array}{l} \dim_ {k} (\mathrm{Hom} _ {\mathrm{A} _ {\mathrm{Q}}} (\mathrm{M}, \mathrm{N})) - \dim_ {k} (\mathrm{Ext} _ {\mathrm{A} _ {\mathrm{Q}}} ^ {1} (\mathrm{M}, \mathrm{N})) \\ = \sum_ {s \in \mathrm{Som} (\mathrm{Q})} \dim (\mathrm{M} _ {s}) \dim (\mathrm{N} _ {s}) - \sum_ {f \in \mathrm{Fl} (\mathrm{Q})} \dim (\mathrm{M} _ {o (f)}) \dim (\mathrm{N} _ {t (f)}). \end{array}
\]

\(d\)) 设 \(k\) 是一个交换域。证明对 Q 的每一对顶点 \((u,v)\)，向量空间 \(\mathrm{Ext}_{\mathrm{A_Q}}^1 (k(u),k(v))\) 的维数等于 Q 中以 \(v\) 为起点、\(u\) 为终点的箭的集合的基数。

\begin{enumerate}
\def\labelenumi{\arabic{enumi})}
\setcounter{enumi}{9}
\item
  设 \(k\) 是一个代数闭域。设 Q 与 Q' 是每个圈的长均为零的有限箭图。证明代数 A \(_{Q'}\) 与代数 A \(_{Q}\) Morita 等价（A, VIII, §6）当且仅当箭图 Q 与 Q' 同构。（取 \(k\)-单模 \(k(s)\)（\(s \in \text{Som}(Q)\)）作为模 M、N。）
\end{enumerate}

